% =============================================================================
% CAPÍTULO II - MARCO TEÓRICO DEL PROYECTO TECNOLÓGICO
% =============================================================================

\cleardoublepage
\thispagestyle{empty}
\vspace*{\fill}
\begin{center}
{\Large\bfseries CAPÍTULO II}\\[2.5cm]
{\Large\bfseries MARCO TEÓRICO DEL PROYECTO TECNOLÓGICO}
\end{center}
\vspace*{\fill}
\clearpage

\chapter*{}
\addcontentsline{toc}{chapter}{CAPÍTULO II -- MARCO TEÓRICO DEL PROYECTO TECNOLÓGICO}
\setcounter{chapter}{2}
\setcounter{section}{0}
\renewcommand{\thesection}{2.\arabic{section}}
\renewcommand{\thesubsection}{2.\arabic{section}.\arabic{subsection}}

\vspace{-3cm}

\section{Marco conceptual}

En el contexto de las evaluaciones por competencia curricular, los métodos de evaluación juegan un papel crucial. Estos métodos son los enfoques y herramientas utilizados para medir el grado en que los estudiantes han adquirido las competencias específicas delineadas en el plan de estudios; pueden incluir desde pruebas tradicionales hasta proyectos prácticos, evaluaciones basadas en la resolución de problemas o presentaciones orales \cite{ref25_gulikers}.

La competencia, en este contexto, se refiere a la capacidad demostrada por un estudiante para aplicar conocimientos, habilidades y actitudes de manera efectiva y apropiada en situaciones específicas. En lugar de simplemente memorizar información, la competencia implica la capacidad de utilizar ese conocimiento de manera significativa y transferible en el mundo real \cite{ref25_gulikers}.

El aprendizaje individualizado es un enfoque educativo que reconoce las diferencias individuales entre los estudiantes y adapta los procesos de enseñanza y evaluación para satisfacer las necesidades únicas de cada uno. En el contexto de las evaluaciones por competencia curricular, esto implica la personalización de las evaluaciones para considerar los estilos de aprendizaje, los intereses y las habilidades individuales de cada estudiante \cite{ref26_tomlinson}.

La retroalimentación significativa se refiere a los comentarios proporcionados a los estudiantes después de una evaluación con el propósito de guiar su aprendizaje de manera efectiva. En el contexto de las evaluaciones por competencia curricular, la retroalimentación significativa se centra en destacar no solo los errores, sino también los logros del estudiante, identificando áreas de mejora y proporcionando orientación concreta sobre cómo avanzar \cite{ref27_blackwiliam}.

La discapacidad intelectual es una condición que afecta la capacidad de una persona para comprender, procesar y retener información, así como para participar en actividades académicas en igualdad de condiciones que sus pares. En el contexto de las evaluaciones por competencia curricular, resulta crucial considerar las necesidades y adaptaciones específicas que puedan requerir los estudiantes con discapacidades intelectuales para participar de manera equitativa en el proceso de evaluación \cite{ref28_aamr}.

El currículo, también conocido como plan de estudios, se refiere al conjunto de objetivos educativos, contenidos, métodos de enseñanza y criterios de evaluación diseñados para guiar el aprendizaje en un programa educativo específico. En el contexto de las evaluaciones por competencia curricular, el currículo proporciona el marco sobre el cual se basan las evaluaciones, delineando las competencias que se espera que los estudiantes adquieran y demuestren durante su proceso educativo \cite{ref29_curriculo_ec}.

\subsection{Evolución de las evaluaciones educativas}

El concepto de evaluación ha evolucionado en consonancia con el concepto de educación predominante. Es así que, desde una evaluación centrada en el acto de juzgar el valor de las cosas, se ha avanzado hacia una evaluación que pretendía asignar valores precisos de medición a determinados objetos educativos \cite{ref30_stufflebeam}.

En la década de los treinta, Ralph Tyler inicia un movimiento de evaluación en función del logro de determinados objetivos formulados con antelación, lo que produjo un cambio relevante para comprender el porqué del proceso educativo, que hasta entonces se enfocaba únicamente en los resultados del aprendizaje \cite{ref31_tyler}.

En los años sesenta la evaluación se redefine como una etapa de \textit{recopilación y uso de información para tomar decisiones sobre un programa educativo}. Este enfoque transforma la percepción previa de la evaluación como un mero instrumento de control y medición para la valoración final de un proceso \cite{ref31_tyler}. En su lugar, se concibe como un medio que proporciona la oportunidad de retroalimentar continuamente el proceso educativo y respaldar la toma de decisiones, mejorando el alcance de los objetivos en el entorno educativo.

A inicios del siglo \textsc{xxi} se tiende a adoptar una concepción ecléctica del proceso de evaluación, definiéndola como el proceso de delinear, obtener, procesar y proporcionar información válida, confiable y oportuna que permita juzgar el mérito o valía de programas, procedimientos y productos con el fin de tomar decisiones \cite{ref32_scriven}.

En la actualidad se considera crucial evaluar la comprensión del estudiante respecto a la actividad a realizar, incluyendo su significado, sentido, plenitud y la manera en que se alcanza esa comprensión, lo que constituye un componente esencial de la orientación necesaria para distinguir cualidades diversas en el proceso de aprendizaje \cite{ref16_jonassen}.

\subsection{Evaluación por competencias}

La evaluación por competencias implica la recopilación de evidencias mediante actividades de aprendizaje y la formulación de valoraciones sobre la medida y la naturaleza del progreso del estudiante en relación con los resultados de aprendizaje esperados \cite{ref33_cano}.

Este modelo evaluativo ofrece nuevas oportunidades a los estudiantes al generar entornos significativos de aprendizaje que acercan sus experiencias académicas al mundo profesional, donde pueden desarrollar capacidades integradas y orientadas a la acción, con el objetivo de resolver problemas prácticos o enfrentarse a situaciones auténticas \cite{ref34_levy}. Estas competencias están compuestas por un conjunto de estructuras de conocimiento, habilidades cognitivas, interactivas y afectivas, actitudes y valores, todas necesarias para la ejecución de tareas, la solución de problemas y el desempeño eficaz en una determinada profesión, organización o rol \cite{ref33_cano,ref34_levy}.

\subsection{Evaluación por competencias en matemáticas}

La competencia en matemáticas constituye un tema de gran relevancia en los países que han adoptado el enfoque competencial para la enseñanza de la materia \cite{ref35_pisa}. Este enfoque destaca la necesidad de reemplazar un currículo centrado únicamente en la adquisición de contenidos para lograr el éxito académico. En su lugar, se aboga por un currículo centrado en la alfabetización matemática de los estudiantes, lo que les permita utilizar de manera integral y efectiva los conocimientos matemáticos en diversas situaciones cotidianas \cite{ref36_niss}. En esencia, se busca enseñar para la vida y no exclusivamente para la escuela.

Desde esta perspectiva, la competencia matemática se entiende como aquella que permite dar un uso funcional al conocimiento matemático en diversas situaciones desde una profunda comprensión \cite{ref35_pisa}. Cuando se enfrentan a problemas en contextos del mundo real, los estudiantes deben activar las competencias matemáticas pertinentes para resolverlos, lo que requiere una cultura matemática que les permita analizar, razonar y comunicar efectivamente al plantear, resolver e interpretar problemas de naturaleza matemática en diversas situaciones \cite{ref37_pisa_marco}.

El desarrollo de la competencia matemática no constituye un proceso de aprendizaje con un punto final \cite{ref34_levy}. Nunca se puede afirmar que tras un determinado proceso educativo se haya adquirido completamente la competencia en cuestión: se trata de un proceso continuo y, cada vez que se enfrentan nuevas situaciones y se relacionan nuevos conocimientos, se incrementa la competencia \cite{ref35_pisa}. Por lo tanto, la evaluación de competencias implica reconocer que el proceso vivido ha influido en el desarrollo de habilidades.

En el marco de un enfoque integrado de enseñanza y aprendizaje, la primera decisión consiste en caracterizar la competencia matemática y sus dimensiones correspondientes \cite{ref36_niss}. Para ello, se deben describir mediante indicadores las acciones que reflejan el progreso en esas dimensiones, configurando un sistema de indicadores que permita describir el dominio alcanzado. Cada acción se define mediante indicadores que muestran las diversas respuestas de los estudiantes según la destreza o habilidad demostrada al resolver la tarea propuesta.

\subsection{Tipos de preguntas en evaluaciones por competencias}

En la evaluación por competencias, el diseño de preguntas juega un papel crucial para medir de manera efectiva las habilidades y conocimientos específicos que se buscan evaluar. Estas preguntas están diseñadas para ir más allá de la simple memorización de información, buscando evidenciar la aplicación práctica de conocimientos y habilidades en situaciones del mundo real \cite{ref38_wiggins}. A continuación, se describen los tipos de preguntas comunes utilizadas en la evaluación por competencias:

\begin{itemize}[leftmargin=*]
\item \textbf{Preguntas de aplicación.} Desafían a los evaluados a aplicar sus conocimientos y habilidades en situaciones prácticas. Por ejemplo: ``Usando tus juguetes, muestra cómo harías un castillo alto con bloques. ¿Puedes explicar qué partes usaste y por qué?''.
\item \textbf{Estudios de caso.} Presentan un escenario detallado y solicitan al evaluado que analice, interprete y proponga soluciones basadas en sus conocimientos. Ejemplo: ``Imagina que tus amigos quieren jugar diferentes juegos. Algunos quieren jugar con carros y otros con muñecas. ¿Cómo podrías ayudarles a decidir qué juego jugar y asegurarte de que todos estén felices?''.
\item \textbf{Problemas de resolución.} Requieren que el evaluado resuelva problemas utilizando sus habilidades y conocimientos, ya sea en situaciones prácticas o teóricas. Por ejemplo: ``Resuelve la siguiente operación y explica cada paso''.
\item \textbf{Preguntas situacionales.} Presentan situaciones hipotéticas o reales y solicitan al evaluado tomar decisiones o proporcionar soluciones basadas en su conocimiento y experiencia. Ejemplo: ``Imagina que estás organizando una fiesta con tus amigos. Algunos quieren jugar dentro de casa y otros fuera. ¿Cómo asegurarías que todos se diviertan y estén contentos?''.
\end{itemize}

Estos tipos de preguntas están diseñados para evaluar no únicamente el conocimiento teórico, sino también la capacidad de aplicar ese conocimiento en contextos prácticos \cite{ref15_villardon}. Adicionalmente, permiten medir habilidades como el pensamiento crítico, la resolución de problemas, la toma de decisiones y la capacidad de actuar en situaciones del mundo real \cite{ref18_brown}.

\subsection{Integración de diferentes tipos de preguntas}

La integración de diversos tipos de preguntas en una evaluación resulta fundamental para capturar de manera integral las habilidades y conocimientos de los estudiantes. Al combinar preguntas de aplicación, estudios de caso, problemas de resolución y preguntas situacionales, se abarcan diferentes niveles de comprensión y se evalúan habilidades específicas, desde la aplicación práctica hasta el pensamiento crítico \cite{ref38_wiggins}. Esta diversidad no solo se adapta a distintos estilos de aprendizaje, sino que también proporciona a los educadores una perspectiva más completa del rendimiento estudiantil, permitiéndoles ofrecer retroalimentación precisa y fomentar el desarrollo integral de los estudiantes.

La importancia radica en la capacidad de estas preguntas para reflejar la complejidad del mundo real y preparar a los estudiantes para enfrentar desafíos diversos. Adicionalmente, motiva la participación activa al mantener la evaluación interesante y relevante. En última instancia, la integración de distintos tipos de preguntas en una evaluación no solo mide el conocimiento superficial, sino que impulsa a los estudiantes a aplicar, analizar y tomar decisiones, promoviendo así un aprendizaje más significativo y orientado al desarrollo de competencias \cite{ref39_haladyna}.

\subsection{Uso de elementos multimedia en preguntas educativas}

La inclusión de elementos multimedia en las evaluaciones dirigidas a niños desempeña un papel crucial al fomentar un ambiente de aprendizaje estimulante y accesible. La incorporación de imágenes coloridas, videos interactivos y elementos visuales no solo capta la atención de los niños, sino que también facilita la comprensión de conceptos y situaciones, lo que promueve una experiencia de evaluación más participativa y emocional \cite{ref40_mayer,ref41_clark}.

Los elementos multimedia, además de resultar atractivos, se alinean con la naturaleza lúdica del aprendizaje infantil. Utilizar imágenes y videos puede convertir la evaluación en una experiencia más amigable y divertida para los niños, contribuyendo a un ambiente de aprendizaje positivo \cite{ref42_cabero}. Asimismo, al simular escenarios de la vida real mediante formatos multimedia, se evalúa la capacidad de los niños para aplicar sus habilidades en contextos prácticos de manera más intuitiva.

La claridad y la comprensión mejoran significativamente al utilizar elementos multimedia adaptados a la edad. Los niños pueden beneficiarse al ver visualmente conceptos abstractos o seguir instrucciones a través de ilustraciones, lo que les permite expresar sus conocimientos y habilidades de manera más efectiva \cite{ref43_paivio}.

\subsection{La tecnología educativa}

La tecnología educativa se ocupa del análisis de los medios, materiales, sitios web y plataformas tecnológicas utilizados en los procesos de aprendizaje. Dentro de su ámbito se incluyen recursos diseñados para abordar las necesidades de los usuarios en términos formativos e instructivos. Las investigaciones del área se centran en el examen del uso de las TIC en entornos de enseñanza y aprendizaje, abordando el impacto general de las tecnologías en la educación a través de enfoques sistémicos, evaluando los procesos mediados desde una perspectiva integral e integradora \cite{ref42_cabero}.

En cuanto al uso de tecnologías en la educación, la conexión entre educación y tecnología no es algo novedoso, sino más bien un elemento constante a lo largo de la historia. En la actualidad se enfatiza que no se trata simplemente de aumentar la intensidad del uso de la tecnología, sino de comprender claramente los beneficios que las opciones tecnológicas pueden aportar \cite{ref44_selwyn}. Resulta especialmente relevante considerar cómo estas alternativas pueden contribuir de manera significativa a que los estudiantes aprendan de forma más efectiva, mejorando y diversificando el proceso educativo.

Actualmente, la tecnología facilita estrategias didácticas enriquecedoras, permitiendo diseños atractivos y motivadores para alcanzar objetivos educativos. Adicionalmente, propicia un cambio de roles tanto para docentes como para estudiantes al aumentar la persistencia, focalización y accesibilidad en las tareas, posibilitando un mayor acompañamiento docente en el proceso de enseñanza-aprendizaje \cite{ref45_unesco_tic}. La tecnología emerge como una herramienta versátil que potencia la interactividad y la colaboración en el entorno educativo.

\subsection{Aplicaciones educativas}

Las aplicaciones educativas representan herramientas valiosas diseñadas para consolidar los conocimientos de los niños en diversas áreas de aprendizaje. Estas aplicaciones, enriquecidas con imágenes, sonidos, dibujos y animaciones adaptadas a las distintas edades y disciplinas, se erigen como elementos cruciales para el proceso educativo \cite{ref46_app_educativas}. Se reconoce que la inclusión de elementos visuales y auditivos facilita el procesamiento de información por parte de los seres humanos.

La relevancia de estas aplicaciones radica en su capacidad para ofrecer un entorno atractivo, estimulante y efectivo, convirtiéndose en factores significativos para potenciar el desarrollo educativo de los niños \cite{ref46_app_educativas}. Resulta imperativo que estas aplicaciones se incorporen progresivamente en diversas instituciones educativas, especialmente considerando su sinergia con internet y el acceso a través de dispositivos móviles, característicos de la era tecnológica actual.

Más allá del aprendizaje individual, estas aplicaciones desempeñan un papel esencial en la promoción de la colaboración y la mejora de la comunicación entre los estudiantes. Facilitan el intercambio de experiencias y el trabajo en equipo, contribuyendo a un ambiente educativo más interactivo y enriquecedor \cite{ref46_app_educativas}.

A pesar de su potencial, la implementación efectiva de estas aplicaciones se ve obstaculizada en parte por la falta de conocimiento y la reticencia de algunos docentes. En la actualidad, aunque surgen constantemente nuevas aplicaciones educativas, son pocos los educadores que las integran en sus clases, lo que podría deberse a la falta de familiaridad o de habilidad para utilizarlas, o a la percepción de que no son adecuadas para el contexto educativo \cite{ref08_ertmer}. Resulta fundamental superar las barreras identificadas y promover la adopción de estas herramientas tecnológicas como aliadas para el desarrollo pedagógico.

\subsection{Desarrollo de aplicaciones para evaluación educativa}

La importancia de las aplicaciones educativas radica en su capacidad para proporcionar experiencias de aprendizaje interactivas y personalizadas. Al incorporar elementos visuales, auditivos y táctiles, estimulan la participación activa de los estudiantes, fomentando un ambiente educativo dinámico y atractivo. Adicionalmente, ofrecen la flexibilidad necesaria para adaptarse a diversos estilos de aprendizaje, permitiendo un enfoque más individualizado y centrado en el estudiante \cite{ref46_app_educativas}.

La selección adecuada de tecnologías para el desarrollo de aplicaciones educativas resulta esencial para garantizar su eficacia y alineación con los objetivos pedagógicos. Resulta crucial que estas tecnologías sean intuitivas, seguras y adaptables a diferentes dispositivos, asegurando un acceso fácil y equitativo para todos los estudiantes \cite{ref48_nielsen}. Adicionalmente, deben integrarse fluidamente en el entorno educativo existente, complementando y enriqueciendo las metodologías tradicionales de enseñanza.

Para el desarrollo de Kiddy Quiz se seleccionaron tecnologías avanzadas que contribuyen significativamente a la eficacia y versatilidad del proyecto. Angular, reconocido por su robustez en la creación de interfaces de usuario interactivas, aporta una experiencia visual envolvente para los estudiantes \cite{ref49_angular}. Su capacidad para gestionar de forma eficiente la interactividad y la presentación de datos en tiempo real se alinea con el objetivo de estimular la participación activa en un entorno educativo dinámico.

Node.js, elegido como entorno de ejecución para el desarrollo del lado del servidor, constituye la columna vertebral que sustenta la aplicación \cite{ref50_nodejs}. Su naturaleza basada en eventos y su capacidad para manejar múltiples conexiones simultáneas garantizan una respuesta rápida y eficiente, proporcionando una experiencia de usuario fluida. La escalabilidad inherente de Node.js asegura que la aplicación pueda adaptarse fácilmente a un número variable de usuarios.

La elección de PostgreSQL como sistema de gestión de bases de datos subraya el compromiso del proyecto con la seguridad, la integridad y la capacidad de adaptación de la aplicación \cite{ref51_postgresql}. PostgreSQL ofrece un sólido soporte para datos estructurados y no estructurados, lo que facilita la gestión de la información educativa de manera eficiente \cite{ref52_postgres_book}. Su capacidad para integrarse de manera transparente con Node.js asegura una comunicación eficaz entre el servidor y la base de datos, optimizando la gestión de datos y proporcionando una base sólida para la adaptabilidad y la personalización del aprendizaje.

\subsection{Programación Extrema (XP)}

La metodología de Programación Extrema (XP), propuesta por Kent Beck \cite{ref53_beck}, surge como un marco de trabajo ágil diseñado para mejorar la calidad del software y la capacidad de respuesta del equipo de desarrollo ante los cambios en los requisitos del cliente. Su enfoque se basa en la sinergia de un conjunto de prácticas de ingeniería y gestión que, aplicadas de manera conjunta y disciplinada, potencian la comunicación, la simplicidad y la retroalimentación continua. XP es independiente de herramientas o lenguajes particulares, brindando a los equipos la flexibilidad para construir software de alta calidad en entornos de alta incertidumbre.

Las prácticas de esta metodología se alinean con un ciclo de vida iterativo y se fundamentan en valores como la comunicación constante, la simplicidad del diseño y el coraje para refactorizar y mejorar el sistema continuamente. El objetivo de XP es entregar valor al cliente a través de lanzamientos frecuentes de software funcional, reduciendo los riesgos y optimizando el proceso de desarrollo mediante la colaboración y la excelencia técnica \cite{ref47_ganney_xp}.

A continuación, se describen las prácticas fundamentales de la Programación Extrema, organizadas según las cuatro actividades principales del ciclo de desarrollo:

\begin{itemize}[leftmargin=*]
\item \textbf{Fase de Planeación.} Esta etapa inicial establece la hoja de ruta del proyecto mediante una colaboración intensa entre el equipo de desarrollo y el cliente. El objetivo principal es definir el alcance del trabajo y priorizar las funcionalidades más valiosas para el negocio. A través del Juego de Planificación, el cliente presenta sus necesidades en forma de Historias de Usuario, las cuales son discutidas, estimadas en esfuerzo por el equipo técnico y priorizadas según su valor. La presencia activa del Cliente en el Sitio resulta fundamental para resolver dudas y tomar decisiones de manera ágil. Adicionalmente, se define una Metáfora, una visión compartida del sistema que facilita la comprensión y la comunicación entre los involucrados.

\item \textbf{Fase de Diseño.} Una vez planificado el trabajo para la siguiente iteración, esta fase se centra en cómo se implementarán las funcionalidades seleccionadas. La prioridad en XP es mantener el Diseño Simple, buscando siempre la solución más sencilla que cumpla con los requisitos actuales. Se evita la complejidad innecesaria y la anticipación de funcionalidades futuras, enfocándose en resolver el problema inmediato de la manera más clara y eficiente posible.

\item \textbf{Fase de Codificación.} La construcción del software se realiza con un fuerte énfasis en la calidad técnica a través de prácticas interrelacionadas. El Desarrollo Dirigido por Pruebas (TDD) es una práctica fundamental, en la que se escriben pruebas automatizadas que fallan antes de escribir el código funcional, guiando así el desarrollo y asegurando que cada parte del sistema cumpla con su propósito. La Refactorización se realiza de forma constante para mejorar la estructura interna del código sin alterar su comportamiento externo. Finalmente, la Integración Continua, la Propiedad Colectiva del Código y los Estándares de Codificación aseguran un código base coherente y de alta calidad.

\item \textbf{Fase de Prueba y Entrega.} En esta etapa el equipo se asegura de que el software construido cumpla con las expectativas del cliente y se entrega valor de forma continua. Los Lanzamientos Pequeños son una práctica clave, en la que se liberan versiones funcionales del sistema con frecuencia, lo que permite obtener retroalimentación temprana que alimenta la planificación de las siguientes iteraciones. La validación se realiza a través de las pruebas unitarias del TDD y, en niveles superiores, mediante pruebas de aceptación que validan la implementación correcta de las historias de usuario.

\item \textbf{Práctica transversal -- Ritmo Sostenible.} Aplicada a lo largo de todas las fases, reconoce la importancia del bienestar del equipo para la productividad a largo plazo. Se promueve un horario de trabajo razonable, con el objetivo de mantener un equipo motivado, enfocado y capaz de producir software de alta calidad de manera constante.
\end{itemize}

\section{Marco referencial}

La evaluación basada en competencias ha experimentado una integración cada vez mayor con la tecnología en los últimos años, lo que ha dado lugar al desarrollo de diversas aplicaciones destinadas a facilitar y mejorar este proceso. A continuación se presentan algunas de las herramientas más representativas del estado del arte:

\textbf{SCALA (Scalable Competence Assessment through a Learning Analytics approach).} Desarrollada en 2015, esta herramienta genera análisis de aprendizaje a partir de datos recopilados de diversos repositorios educativos. SCALA elabora rúbricas detalladas considerando el modelo curricular basado en competencias de la Universidad de Deusto y los datos extraídos de los repositorios sobre las actividades realizadas por los estudiantes durante los cursos, así como sus resultados de evaluación. Estas actividades miden indicadores de éxito relacionados con determinadas competencias y los resultados se presentan en un panel de control \cite{ref21_scala_universidad}.

\textbf{CAT (Competence Assessment Tool).} Plataforma de 2014--2016 que permite la evaluación de competencias genéricas a nivel universitario utilizando blogs. Su objetivo es mejorar los procesos de autorreflexión y proporcionar retroalimentación junto con rúbricas para evaluar indicadores de éxito asociados a competencias específicas. Los resultados de las pruebas realizadas muestran un alto nivel de satisfacción entre todos los involucrados y destacan que la participación de los estudiantes en el proceso de evaluación incrementa su aprendizaje mediante la autorregulación \cite{ref10_cat}.

\textbf{SECAT (Software Engineering Competence Assessment Tool).} Marco de evaluación por competencias para la ingeniería de software desarrollado en 2014--2015, que adapta y enriquece el enfoque de Rauner. Utiliza evaluaciones individuales y de equipo a través de cuestionarios cuyas preguntas están relacionadas con diversas dimensiones de competencia. Se emplea un gráfico radial para describir el nivel alcanzado por el estudiante en cada competencia \cite{ref11_secat,ref22_secat_engineering}.

\textbf{LACAMOLC.} Propuesta de plataforma web de 2014 que ofrece a docentes y estudiantes un panel de control que recopila datos sociales y de uso de diversas tecnologías de aprendizaje, como Moodle y Google Apps \cite{ref23_lacamolc}. Esta plataforma demanda nuevas competencias a los profesionales, que deben poseer habilidades y destrezas específicas. Utiliza Pentaho como herramienta de análisis, lo que permite escalar efectivamente los sistemas y proporcionar visualizaciones para respaldar el proceso de aprendizaje y evaluación.

\textbf{AEEA (Adaptive Engine for Educational Analytics).} Suite de herramientas web adaptativas de 2013 que permite diseñar, evaluar y desplegar análisis de aprendizaje para cursos en línea y mixtos de instituciones de educación superior bajo el marco europeo de cualificaciones (EQF). Los instrumentos de evaluación se basan en formularios de rúbricas para la evaluación integral del estudiante, y las analíticas presentan al estudiante frente a diferentes contextos sociales (individual, trabajo en grupo, curso) y su rendimiento comparativo en las diferentes competencias evaluadas \cite{ref24_aeea}.

Estas aplicaciones han surgido como soluciones que abordan diversos aspectos clave, tales como el análisis de datos mediante paneles de control y gráficos intuitivos, la capacidad de adaptarse a una variedad de contextos para los cuestionarios, la personalización de las evaluaciones y la alineación con los objetivos del plan de estudios. Sin embargo, ninguna de ellas integra todas estas capacidades de manera unificada, además de que suelen omitir aspectos importantes como la variabilidad de tipos de preguntas o la creación de cuestionarios dentro de la propia aplicación. Resulta especialmente relevante señalar que todas estas herramientas fueron diseñadas para el ámbito de la educación superior. Es en este contexto donde Kiddy Quiz emerge como una solución diferenciadora, con el propósito de abordar todas estas áreas de manera integral, centrándose particularmente en los estudiantes de educación básica elemental y aprovechando los recientes avances en inteligencia artificial generativa para producir retroalimentación adaptativa \cite{ref46_app_educativas}.
